\section{Teorema de Bayes}
\begin{thm:bayes}[Teorema de Bayes]
Seja $(\Omega, \mathcal{B}, P)$ um espaço de probabilidade.
Se $E$ e $F$ são eventos, então
\[
	P(E | F) = P(F | E) \frac{P(E)}{P(F)}
\]
\begin{proof}
	\begin{gather*}
		P(E | F) \cdot P(F) = P(E \cap F) = P(F | E) P(E)\\
		P(E | F) = P(F | E) \frac{P(E)}{P(F)}
	\end{gather*}
\end{proof}
\end{thm:bayes}

Uma interpretação para o teorema é a atualização na crença $P(E)$, chamada de
\textbf{prior}, numa hipótese $E$ dada uma nova evidência $F$ para ela, obtendo um novo
grau de crença $P(E|F)$ chamado de \textbf{posterior}. A probabilidade $P(F|E)$ é
chamada de \textbf{verossimilhança}, denotando o quão verossímil é observar a evidência $F$
considerando a hipótese $E$ verdadeira. Por sua vez, $P(F)$ é a crença na ocorrência
de $F$ tendo em vista todas as hipóteses, que, junto com a verossimilhança, ajusta
o novo grau de crença em $E$.

Se num dado experimento aleatório consideramos hipóteses $E_1, \ldots, E_n$ e obtemos
uma evidência $F$, é plausível questionar em qual hipótese deveríamos depositar maior
crença visto a evidência. Conseideremos então o corolário:

\begin{cor:formula de bayes}[Fórmula de Bayes]
Seja $\mathcal{E} = \{E_i\}_\mathbb{N}$ uma partição de eventos de um espaço amostral
$\Omega$.
Se $F$ é um evento de $\Omega$, então
\[
	P(E_i | F) = \frac{P(F | E_i) P(E_i) }{\sum_{E \in \mathcal{E}} P(E) P(F | E)}
\]
\begin{proof}
	Segue do Teorema de Bayes e do Teorema da Probabilidade Total.
\end{proof}
\end{cor:formula de bayes}

Logo, para cada hipótese, poderíamos utilizar a Fórmula de Bayes para computar a crença
\textit{a posteriori} nela após a evidência $F$, e selecionar aquela que tivesse maior
probabilidade.
